\section{Evolution of experimental methodology and integrity controls}
\label{sec:history}

PHASE-TTC developed from predictive event representations into geometric motion estimation, phase-domain temporal refinement, and event--RGB fusion. The methodological changes below explain how evaluation controls evolved alongside the model. Historical experiment identifiers and configurations remain in the evidence registry; Table~\ref{tab:source-revisions} collects the associated code revisions.

\subsection{Predictive representations and causal measurements}
Initial experiments established indexed EvTTC access, voxel representations, a supervised TinyCNN, and geometric and event-rate references. Multihorizon predictive learning, frozen-encoder probes, and low-label comparisons tested the original E-JEPA-TTC hypothesis. Token transformers, tubelets, navigation conditioning, and rollout probes extended the temporal models.

Causal validity required changes to both the measurements and evaluation procedure. Centered geometric derivatives used future observations and were replaced with past-only derivatives. This annotation-assisted reference was evaluated on its supported queries, with coverage reported separately from methods using different inputs. The same distinction later informed the separation of sensor support, prediction availability, and complete-system accuracy.

\subsection{Sequence isolation and initialization dependence}
The early CPLA-high comparisons reused inspected sequences and shared one self-supervised initialization across three fine-tuning seeds. They consequently measure downstream optimization variability on development data. Checkpoint provenance, cache versions, nested label subsets, and matched-scratch comparisons were added to the protocol, and these runs were excluded from independent-test conclusions.

A separate validation-path error recorded completion without executing the intended evaluation matrix. The corrected path executes the matrix and checks its results; the earlier completion markers provide no experimental evidence. This correction distinguishes verified evaluation outputs from the status of a software run.

The first grouped geometry study also initialized held-out folds from a representation trained on all nine development sequences. Fold-specific producer training removed that overlap before downstream fitting. Checkpoint selection and earlier sequence exposure still define these studies as development comparisons. \evidence{E03}

\subsection{Geometry, transport, and controlled model selection}
Object-centred experiments separated absolute geometry from changes in apparent scale. Foreground supervision, RGB-derived segmentation masks, and relational distillation tested spatial information; cross-time transport tested motion correspondence. Their supervision sources are recorded separately from their inference modalities.

Scaling the relationally distilled student from 2,048 to 8,192 training examples while increasing its train-selected distillation weight reduced MiD from 322.7 to 262.8 at unchanged capacity. The missing same-size $\lambda=4$ control prevents separate attribution to data and loss weight. This single-seed result preceded the fold-local recovery studies. \evidence{H_A4_S1}

The fold-local transport studies compared accuracy and geometric retention under matched samples and optimization budgets. No model in the two initial recovery campaigns met all selection criteria. Complementary errors between the baseline and coarse-to-fine producers then motivated a prospective nested-routing comparison, accompanied by frozen, partial-finetuning, random, scratch, and shuffled-future predictive-learning controls. \evidence{E03--E11,E16--E24}

\subsection{Temporal context and multimodal estimation}
Phase-domain studies examined incremental transport information, raw timing, expert routing, risk estimation, and context with fixed optimization budgets. The original residual-context campaign completed 72 fits and 180,000 saved updates. Its eight-slot history improved on the current-observation head by 11.3 MiD across three head seeds, with a hierarchical 95\% difference interval of $[-16.2,-6.6]$ on reused development data. Trained reversal, repeated-current, selection, and latent-input controls constrain the mechanism interpretation (Table~\ref{tab:original-t6}). A later extension completed 24 fits and 60,000 updates, reusing original H8 endpoints; the two totals describe distinct campaigns. Short- and long-history heads shared upstream producers. Capacity comparisons retained incomplete training endpoints as partial records. \evidence{H_T6,M_SIMPLEX}

Efficient input preparation and the RGB port added native sensor clocks and transfer evaluation.

The intervening multimodal selection campaign completed 32 fits and 80,000 updates. Its phase-only correction passed the original 2\% band safeguards in three head seeds, retaining seed 7. This historical selection used in-sample stacking on H; a separate four-fold H cross-fitting experiment failed the MAE criterion. Both differ from the subsequent five-rotation TRAIN40 refit. Section~\ref{sec:v16-selection} records their roles and the different event controls used in V16 and V18. \evidence{H_V16}

Later range and representation studies tested clipping, nonlinear readouts, geometry bypasses, encoder adaptation, and loss weighting. Several variants improved individual TTC bands while worsening receding-motion estimates or another selection metric. None of the joint-measurement variants met the complete selection criteria. The earlier phase-correction recipe remained the reference for the final refit, while these unsuccessful variants motivated separate negative-range controls. \evidence{M_SIMPLEX,L02,L10,L24,L25,L30}

\subsection{Final refitting, adaptation, and execution}
The submitted TRAIN40 lineage follows the phase-correction recipe, using group-separated auxiliary models for out-of-fold fitting and all forty sequences for deployed producers. The official GarlTTC Test evaluation covers the resulting complete system (Section~\ref{sec:official}).

Runtime studies examined persistent histories, exact normalization, graph execution, crop reuse, and multiple objects. Timing separates initialization from measured arrivals and successful outputs from failed requests. Smaller event packets preserved numerical parity in the tested setting; throughput requires a separate paced measurement.

Subsequent work recovered all 434 local EvTTC targets, fitted the quality-conditioned adapter, and tested canonical cadence throughout the frozen estimator. The adapter and later cadence refits remain separate variants. Predictive-learning and physical-measurement controls, together with a separate detection study, examine the remaining representation and localization dependencies. Sections~\ref{sec:quality-adaptation} and~\ref{sec:canonical-cadence}, and Appendices~\ref{app:physics-followup} and~\ref{app:detection-followup}, report their protocols and outcomes.
